\documentclass[11pt]{article}
\usepackage[margin=2.6cm]{geometry}
\usepackage{amsmath,amssymb,amsthm,mathtools}
\usepackage{booktabs}
\usepackage{enumitem}
\usepackage[hidelinks]{hyperref}

\newtheorem{theorem}{Theorem}[section]
\newtheorem{lemma}[theorem]{Lemma}
\newtheorem{proposition}[theorem]{Proposition}
\newtheorem{corollary}[theorem]{Corollary}
\theoremstyle{definition}
\newtheorem{definition}[theorem]{Definition}
\newtheorem{remark}[theorem]{Remark}
\newtheorem{problem}[theorem]{Open problem}

\newcommand{\F}{\mathbb{F}}
\newcommand{\K}{\mathbb{K}}
\newcommand{\ML}{\mathrm{ML}}
\newcommand{\Enc}{\mathrm{Enc}}
\newcommand{\fold}{\mathrm{fold}}
\newcommand{\dpos}{d_{\mathrm{pos}}}
\newcommand{\dblk}{d_{\mathrm{blk}}}
\newcommand{\GL}{\mathrm{GL}}
\newcommand{\Gs}{\mathcal{G}^{*}}
\newcommand{\eq}{\mathrm{eq}}

\title{SPARK: Matrix-Gate Folding\\[2pt]
\large A multilinear polynomial commitment IOP}
\author{Abdelali Mkhida \\ Algorizk}
\date{October 2026}

\begin{document}
\maketitle

\begin{abstract}
We construct a family of linear codes for multilinear polynomials from $\GL_2$ gate matrices
attached to the nodes of a binary tree, a folding operator that is exactly the gate rule
evaluated at a random point, and a commitment scheme in which the folding challenges are the
opening point. The construction uses no structured evaluation domain: gates are arbitrary
elements of $\GL_2(\F)$ with nonzero second row, and the encoder uses only operations of the
form $u+t\,v$. We prove the algebraic identities; three distance theorems for randomly sampled per-node gates
(a one-layer bound, a rank-saturation bound, and a bound with a checked setup); that the generic
code is MDS; a zero decomposition and an obstruction showing what any sharper distance bound must
pay; a proximity lemma and an agreement lemma for the folding operator; and soundness of the
commitment with a per-query detection rate of $\min(1-(1-\Delta)^{1/3},\Delta/2)$, with no loss
in the number of variables. The construction works over any field large enough for the distance
bounds; gates may be sampled in an extension field. A Rust implementation accompanies this
document.
\end{abstract}

\tableofcontents

%=====================================================================
\section{Notation}
%=====================================================================

\begin{itemize}[leftmargin=*]
\item $\F$ is a finite field with $q=|\F|$ (any characteristic, including $2$); $\K\supseteq\F$ is an
extension field used for verifier challenges.
\item $\ML(x_j,\dots,x_n)$ denotes multilinear polynomials in the listed variables;
$\ML_i$ denotes multilinear polynomials in $i$ variables.
\item $k\ge 0$ is the \emph{redundancy parameter}. For $0\le i\le n$ put
\[
\Omega_i=\{0,1\}^k\times\{0,1\}^{i},\qquad L_i=|\Omega_i|=2^{k+i},
\]
and identify $\Omega_i=\Omega_{i-1}\times\{0,1\}$, writing $\omega=(\omega',b)$.
\item A position $\omega\in\Omega_n$ is written $\omega=(\beta,\theta_n,\theta_{n-1},\dots,\theta_1)$
with $\beta\in\{0,1\}^k$, and its prefixes are
$\omega^{(j)}=(\beta,\theta_n,\dots,\theta_{j+1})\in\Omega_{n-j}$ for $0\le j\le n$.
\item For $w,w'\in\K^{\Omega_i}$:
$\dpos(w,w')=L_i^{-1}\,|\{\omega:w(\omega)\neq w'(\omega)\}|$, and for $i\ge1$ the
\emph{block distance}
\[
\dblk(w,w')=L_{i-1}^{-1}\,\bigl|\{\omega'\in\Omega_{i-1}:\ (w(\omega',0),w(\omega',1))\neq(w'(\omega',0),w'(\omega',1))\}\bigr|.
\]
Always $\dpos\le\dblk\le 2\dpos$. Distances to a set are minima.
For a linear code $C$, $\Delta(C)$ is its minimum relative distance.
\end{itemize}

%=====================================================================
\section{Gates}
%=====================================================================

\begin{definition}[Gate]
A \emph{gate} is a matrix
$A=\begin{pmatrix}a&b\\c&d\end{pmatrix}\in\GL_2(\F)$ with $c\neq0$, $d\neq0$.
Write $\delta=ad-bc$ and $\Gs=\{A\in\GL_2(\F):c\neq0,\ d\neq0\}$.
Its \emph{gate weights} are $W_L(x)=a+cx$ and $W_R(x)=b+dx$, and its \emph{structural points} are
\[
t^0=-\frac{a}{c},\qquad t^1=-\frac{b}{d}.
\]
\end{definition}

\begin{lemma}[Gates are point pairs with scales]\label{lem:gate-param}
The map $(c,d,t^0,t^1)\mapsto\begin{pmatrix}-t^0c&-t^1d\\c&d\end{pmatrix}$ is a bijection from
$(\F^\times)^2\times\{(t^0,t^1)\in\F^2:t^0\neq t^1\}$ onto $\Gs$, and
$\delta=cd\,(t^1-t^0)$. In particular, if $A$ is uniform on $\Gs$ then $(t^0,t^1)$ is uniform on
pairs of distinct elements.
\end{lemma}
\begin{proof}
With $a=-t^0c$, $b=-t^1d$ one gets $ad-bc=-t^0cd+t^1cd=cd(t^1-t^0)$, nonzero iff $t^0\neq t^1$.
The inverse map is $A\mapsto(c,d,-a/c,-b/d)$. The fibre over each $(t^0,t^1)$ has size $(q-1)^2$.
\end{proof}

\begin{lemma}[Decomposition and collapse]\label{lem:collapse}
Let $A\in\Gs$ and let $F=u+x\,v$ with $u,v$ independent of $x$. There are unique $F_L,F_R$
independent of $x$ with
\[
F=W_L(x)\,F_L+W_R(x)\,F_R,\qquad F_L=\delta^{-1}(du-bv),\quad F_R=\delta^{-1}(-cu+av).
\]
Moreover
\[
F\big|_{x=t^0}=\gamma^0F_R,\quad \gamma^0=-\frac{\delta}{c},\qquad\qquad
F\big|_{x=t^1}=\gamma^1F_L,\quad \gamma^1=+\frac{\delta}{d}.
\]
\end{lemma}
\begin{proof}
$W_LF_L+W_RF_R=(aF_L+bF_R)+x(cF_L+dF_R)$, so the identity is $A\binom{F_L}{F_R}=\binom{u}{v}$,
uniquely solvable since $\delta\neq0$. At $x=t^0$: $W_L(t^0)=0$ and
$W_R(t^0)=b-ad/c=-\delta/c$. At $x=t^1$: $W_R(t^1)=0$ and $W_L(t^1)=a-bc/d=\delta/d$.
\end{proof}

\begin{lemma}[Gate rule at an arbitrary point]\label{lem:gate-rule}
In the setting of Lemma~\ref{lem:collapse}, put $w^0=F|_{x=t^0}$ and $w^1=F|_{x=t^1}$. For every $z$,
\[
F\big|_{x=z}=W_L(z)\,\frac{d}{\delta}\,w^1-W_R(z)\,\frac{c}{\delta}\,w^0
= w^0+\frac{z-t^0}{t^1-t^0}\,(w^1-w^0).
\]
\end{lemma}
\begin{proof}
By Lemma~\ref{lem:collapse}, $F_L=(d/\delta)w^1$ and $F_R=-(c/\delta)w^0$; substitute into
$F=W_LF_L+W_RF_R$ at $x=z$. The second expression is the unique affine function of $z$ taking the
values $w^0,w^1$ at $t^0,t^1$.
\end{proof}

\begin{lemma}[One-layer sum rule]\label{lem:sumrule}
If $F=W_L(x)F_L+W_R(x)F_R$ then
$\sum_{x\in\{0,1\}}F=(2a+c)\,F_L+(2b+d)\,F_R$.
\end{lemma}
\begin{proof}
$\sum_{x\in\{0,1\}}(a+cx)=2a+c$ and $\sum_{x\in\{0,1\}}(b+dx)=2b+d$.
\end{proof}

%=====================================================================
\section{The gate code}\label{sec:code}
%=====================================================================

\begin{definition}[Gate family]
A \emph{gate family} is a collection
$\mathcal{A}=\bigl(A_j(\omega')\bigr)_{1\le j\le n,\ \omega'\in\Omega_{n-j}}$ of gates, one per
variable $x_j$ and per node $\omega'\in\Omega_{n-j}$. Write $t_j^b(\omega')$ for the structural
points of $A_j(\omega')$. The family is \emph{per-layer} if $A_j(\omega')$ does not depend on
$\omega'$, and \emph{per-node} otherwise.
\end{definition}

\begin{definition}[Encoding]\label{def:enc}
For $0\le i\le n$ let $C_i$ encode $\ML(x_{n-i+1},\dots,x_n)$ as follows.
\[
\Enc_0(c)(\beta)=c\quad(\beta\in\{0,1\}^k),
\]
and for $i\ge1$, writing $j=n-i+1$ and $F=u+x_j\,v$ with $u,v\in\ML(x_{j+1},\dots,x_n)$,
\[
\Enc_i(F)(\omega',b)=\Enc_{i-1}(u)(\omega')+t_j^b(\omega')\,\Enc_{i-1}(v)(\omega'),
\qquad \omega'\in\Omega_{i-1},\ b\in\{0,1\}.
\]
Put $C_i=\Enc_i(\ML(x_{n-i+1},\dots,x_n))\subseteq\F^{\Omega_i}$.
\end{definition}

\begin{proposition}[Closed form]\label{prop:closed}
For $F\in\ML(x_1,\dots,x_n)$ and $\omega=(\beta,\theta_n,\dots,\theta_1)\in\Omega_n$,
\[
\Enc_n(F)(\omega)=F\bigl(r(\omega)\bigr),\qquad r_j(\omega)=t_j^{\theta_j}\bigl(\omega^{(j)}\bigr).
\]
The coordinate $r_j(\omega)$ depends only on $(\beta,\theta_n,\dots,\theta_j)$. Hence two positions
that differ only in $\theta_j$ have points differing only in the coordinate $x_j$.
\end{proposition}
\begin{proof}
Induction on $n$. For $F=u+x_1v$, $\Enc_n(F)(\omega^{(1)},\theta_1)=u(r_{\ge2})+t_1^{\theta_1}(\omega^{(1)})\,v(r_{\ge2})=F(r(\omega))$,
where $r_{\ge2}$ is the point attached to $\omega^{(1)}$ by the induction hypothesis.
\end{proof}

\begin{remark}[Relation to the bottom table]
By Lemma~\ref{lem:collapse}, the bottom table $T_0$ of the original symbolic folding protocol satisfies
$T_0[\theta]=\Gamma(\theta)^{-1}F(r(\theta))$ with a public nonzero $\Gamma(\theta)$. So $T_0$ is
$\Enc_n(F)$ for $k=0$ and a per-layer family, up to a public diagonal scaling, which preserves all
distances. Propositions~\ref{prop:perlayer} and~\ref{prop:rate1} explain why that table alone
cannot support spot-checks.
\end{remark}

\begin{proposition}[Encoding cost]\label{prop:cost}
$\Enc_n(F)$ is computed with at most $(n+1)\,2^{n+k}$ multiplications and as many additions in $\F$.
\end{proposition}
\begin{proof}
With $T(i)$ the cost for $i$ variables, $T(i)=2T(i-1)+L_i$ and $T(0)=L_0=2^k$, so
$T(n)=(n+1)2^{n+k}$.
\end{proof}

\begin{lemma}[Extension of scalars]\label{lem:ext}
For a linear code $C\subseteq\F^N$, the code $C\otimes\K\subseteq\K^N$ has $\Delta(C\otimes\K)=\Delta(C)$.
\end{lemma}
\begin{proof}
Fix an $\F$-basis $(e_l)$ of $\K$. A word of $C\otimes\K$ is $\sum_le_lc_l$ with $c_l\in C$. At a
position $p$, $\sum_le_lc_l(p)=0$ iff all $c_l(p)=0$, so the support is $\bigcup_l\operatorname{supp}c_l$.
\end{proof}

\begin{proposition}[Per-layer families have distance $2^{-n}$]\label{prop:perlayer}
If $\mathcal{A}$ is per-layer, then $\Delta(C_n)\le 2^{-n}$.
\end{proposition}
\begin{proof}
All points $r(\omega)$ lie in the grid $\prod_j\{t_j^0,t_j^1\}$, each grid point appearing $2^k$
times. $F=\prod_{j=1}^n(x_j-t_j^0)$ is nonzero exactly at the grid point
$(t_1^1,\dots,t_n^1)$, so $\operatorname{wt}(\Enc_n(F))=2^k$, of relative weight $2^{-n}$.
\end{proof}

\begin{proposition}[No redundancy, no testing]\label{prop:rate1}
If $k=0$ and $\Delta(C_n)>0$, then $C_n=\F^{\Omega_n}$. Every table is then a codeword, and
$\Delta(C_n)=2^{-n}$.
\end{proposition}
\begin{proof}
$\Enc_n$ is injective and $\dim C_n=2^n=L_n$.
\end{proof}

\subsection{Distance of per-node random families}

\begin{lemma}[One-layer distance]\label{lem:layer}
Let $C\subseteq\F^{N}$ be linear with $\dim C=D$ and $\Delta(C)\ge\Delta$. Let
$(t^0_p,t^1_p)_{p\le N}$ be independent, each uniform on pairs of distinct elements of $\F$, and
\[
C'=\bigl\{(c_u+t^0\circ c_v,\ c_u+t^1\circ c_v)\ :\ c_u,c_v\in C\bigr\}\subseteq\F^{2N},
\]
with coordinates paired as $(p,0),(p,1)$. For every integer $s\ge0$,
\[
\Pr\Bigl[\Delta(C')<\Delta-\tfrac{s}{2N}\ \text{ or }\ \dim C'<2D\Bigr]\ \le\ q^{2D}\binom{N}{s}\Bigl(\frac{2}{q}\Bigr)^{s}.
\]
\end{lemma}
\begin{proof}
Fix $(c_u,c_v)\neq(0,0)$ and let $w$ be the corresponding word of $C'$.

If $c_v=0$, then $w=(c_u,c_u)$ and $\operatorname{wt}(w)=2\operatorname{wt}(c_u)\ge2\Delta N$.

If $c_v\neq0$, consider $p\in\operatorname{supp}c_v$. The pair
$(c_u(p)+t^0_pc_v(p),\ c_u(p)+t^1_pc_v(p))$ has at most one zero, since $t^0_p\neq t^1_p$.
It has a zero iff $\alpha_p:=-c_u(p)/c_v(p)\in\{t^0_p,t^1_p\}$, an event of probability exactly
$2/q$, independently over $p$. Let $Z$ be the set of such $p$. Then
$\operatorname{wt}(w)\ge 2\operatorname{wt}(c_v)-|Z|\ge2\Delta N-|Z|$, and
$\Pr[|Z|\ge s]\le\binom{N}{s}(2/q)^s$.

A union bound over the $q^{2D}$ pairs gives: outside the stated event every nonzero pair has
relative weight $\ge\Delta-s/(2N)>0$ (when this is positive), which also gives injectivity, hence $\dim C'=2D$.
\end{proof}

\begin{theorem}[Distance of the gate code]\label{thm:distance}
Let $\mathcal{A}$ be a per-node family with independent gates, each uniform on $\Gs$. For integers
$s_1,\dots,s_n\ge0$, with probability at least
\[
1-\sum_{i=1}^{n}q^{2^{i}}\binom{L_{i-1}}{s_i}\Bigl(\frac{2}{q}\Bigr)^{s_i},
\]
every $C_i$ $(0\le i\le n)$ has dimension $2^i$ and
\[
\Delta(C_i)\ \ge\ \Delta^{*}:=1-\sum_{l=1}^{n}\frac{s_l}{2L_{l-1}}.
\]
\end{theorem}
\begin{proof}
$C_0$ is the repetition code, with $\Delta(C_0)=1$ and dimension $1$. Definition~\ref{def:enc}
builds $C_i$ from $C_{i-1}$ exactly as $C'$ from $C$ in Lemma~\ref{lem:layer}, with
$N=L_{i-1}$, $D=2^{i-1}$ and fresh structural points (uniform distinct pairs by
Lemma~\ref{lem:gate-param}). Apply Lemma~\ref{lem:layer} at each level and take a union bound.
\end{proof}

\begin{remark}[Parameters]
Using $\binom{N}{s}\le(eN/s)^s$, Table~\ref{tab:bound} lists the smallest certified $\Delta^{*}$
with setup failure probability $\le2^{-128}$. Sample gates in the base field $\F$. By
Lemma~\ref{lem:ext}, the distance carries over to $\K$.
\end{remark}

\begin{table}[h]
\centering
\begin{tabular}{rrrr}
\toprule
$\log_2q$ & $n$ & $k$ & certified $\Delta^{*}$\\
\midrule
31 & 16 & 5 & 0.212\\
31 & 16 & 6 & 0.589\\
31 & 20 & 6 & 0.508\\
31 & 24 & 6 & 0.427\\
62 & 20 & 6 & 0.604\\
124 & 20 & 5 & 0.287\\
124 & 20 & 6 & 0.640\\
\bottomrule
\end{tabular}
\caption{Certified distance (Theorem~\ref{thm:distance}), setup failure $\le 2^{-128}$.
For $k\le4$ the certificate is vacuous at these $n$.}\label{tab:bound}
\end{table}

\begin{table}[h]
\centering
\begin{tabular}{rrrrrrr}
\toprule
$q$ & $n$ & $k$ & length & per-node min & per-node mean & per-layer\\
\midrule
13 & 2 & 3 & 32 & 0.563 & 0.625 & 0.250\\
61 & 2 & 3 & 32 & 0.781 & 0.797 & --\\
7 & 3 & 2 & 32 & 0.250 & 0.339 & 0.125\\
7 & 3 & 3 & 64 & 0.375 & 0.412 & 0.125\\
3 & 4 & 2 & 64 & 0.078 & 0.096 & 0.0625\\
\bottomrule
\end{tabular}
\caption{Exhaustive minimum relative distance over random families (6 samples, 4 for $q=61$).
Per-layer families sit exactly at $2^{-n}$ (Proposition~\ref{prop:perlayer}). Per-node
distance grows with $q$, as Lemma~\ref{lem:layer} predicts ($2/q$ collision rate).}
\label{tab:exp}
\end{table}

\subsection{Distance by rank saturation}

Theorem~\ref{thm:distance} loses about $2^{-k}$ per level. The argument below loses only about
$H_2(2^{-k})/\log_2q$ per level, so large fields allow small redundancy.

For a set $X$ of positions of $C_i$ write $V_X=\{F:\Enc_i(F)|_X=0\}$ and $d_i(X)=\dim V_X$.

\begin{lemma}[Augmentation]\label{lem:augment}
Let $C=\Enc(\ML)$ be injective and let $Z\ge1$. If $V_S=0$ for every $S$ with $|S|\ge Z$, then
$d(X)\le\max(0,Z-|X|)$ for every $X$.
\end{lemma}
\begin{proof}
Suppose $d(X)=d\ge1$. Since $\Enc$ is injective, some element of $V_X$ is nonzero at some
position $\omega$, and $d(X\cup\{\omega\})=d-1$. Repeating $d-1$ times gives $Y$ with $|Y|=d-1$
and $V_{X\cup Y}\neq0$. Hence $|X|+d-1<Z$.
\end{proof}

\begin{lemma}[One level]\label{lem:onelevel}
Fix the gates of $C_{i-1}$, and assume $d_{i-1}(X)\le\max(0,Z'-|X|)$ for all $X$. Let
$S\subseteq\Omega_i$ with $|S|\ge2Z'$. Over the level-$i$ gates,
\[
\Pr\bigl[V_S\neq0\bigr]\ \le\ \frac{(1+1/q)^{L_{i-1}}}{q-1}\;q^{\,2Z'-|S|}.
\]
\end{lemma}
\begin{proof}
Let $T_2$ be the blocks with both entries in $S$ and $T_1$ the blocks with exactly one, say
$(\omega',b_{\omega'})$; then $|S|=2|T_2|+|T_1|$. Write $F=u+xv$ and $U=V_{T_2}$ (level $i-1$),
and put $\hat u=\Enc_{i-1}(u)$, $\hat v=\Enc_{i-1}(v)$.

\emph{Constraints.} $F\in V_S$ iff:
\begin{itemize}[leftmargin=*]
\item $u,v\in U$: on a double block, $\hat u+t^0\hat v=\hat u+t^1\hat v=0$ with $t^0\neq t^1$ forces
$\hat u=\hat v=0$;
\item for every $\omega'\in T_1$: $\hat u(\omega')+t\,\hat v(\omega')=0$, where
$t=t^{b_{\omega'}}(\omega')$ is uniform on $\F$ and independent across blocks.
\end{itemize}
For fixed $(u,v)$, the constraint at $\omega'\in T_1$ holds with probability $1/q$ if
$\hat v(\omega')\neq0$, and with probability $[\hat u(\omega')=0]$ if $\hat v(\omega')=0$.

\emph{Counting.} $V_S$ is a subspace, so $\Pr[V_S\neq0]\le\mathbb{E}[\#(V_S\setminus0)]/(q-1)$.
\begin{itemize}[leftmargin=*]
\item Pairs with $v=0$ contribute $q^{d_{i-1}(T_2\cup T_1)}-1$. This is $0$, because
$|T_2|+|T_1|\ge|S|/2\ge Z'$.
\item For $v\neq0$, let $N(v)=\{\omega'\in T_1:\hat v(\omega')=0\}$. Summing over $u$ gives the weight
$q^{|N(v)|-|T_1|}\,q^{d_{i-1}(T_2\cup N(v))}$.
\item Group the $v\neq0$ by $N=N(v)$. They lie in $V_{T_2\cup N}\setminus0$, which is nonempty only
if $d_{i-1}(T_2\cup N)\ge1$, i.e.\ $|T_2|+|N|<Z'$. Then $d_{i-1}(T_2\cup N)\le Z'-|T_2|-|N|$.
\end{itemize}
Hence
\[
\mathbb{E}[\#(V_S\setminus0)]\le\sum_{N\subseteq T_1}q^{\,|N|-|T_1|+2(Z'-|T_2|-|N|)}
=q^{\,2Z'-|S|}\sum_{N\subseteq T_1}q^{-|N|}=q^{\,2Z'-|S|}(1+1/q)^{|T_1|}.\qedhere
\]
\end{proof}

\begin{theorem}[Distance by rank saturation]\label{thm:distance2}
Let the gates be independent and uniform on $\Gs(\F)$. Choose integers $g_1,\dots,g_n\ge0$, put
$Z_0=1$ and $Z_i=2Z_{i-1}+g_i$, and assume $Z_i\le L_i$. With probability at least
\[
1-\sum_{i=1}^{n}\binom{L_i}{Z_i}\frac{(1+1/q)^{L_{i-1}}}{q-1}\,q^{-g_i},
\]
every nonzero codeword of every $C_i$ has fewer than $Z_i$ zeros. In particular
$\Delta(C_i)\ge1-(Z_i-1)/L_i$, and $\min_i\Delta(C_i)\ge1-(Z_n-1)/L_n$.
\end{theorem}
\begin{proof}
Induction on $i$.
\begin{itemize}[leftmargin=*]
\item \emph{Base.} $C_0$ is the repetition code, so $V_S=0$ for every nonempty $S$, i.e.\ $Z_0=1$.
\item \emph{Step.} Assume the statement for $C_{i-1}$. Then $\Enc_{i-1}$ is injective (because
$Z_{i-1}\le L_{i-1}$), and Lemma~\ref{lem:augment} gives $d_{i-1}(X)\le\max(0,Z_{i-1}-|X|)$ for all $X$.
\item \emph{Union bound.} By Lemma~\ref{lem:onelevel} with $Z'=Z_{i-1}$, each $S$ with $|S|=Z_i$
satisfies $V_S\neq0$ with probability at most $\frac{(1+1/q)^{L_{i-1}}}{q-1}q^{-g_i}$. A union bound
over the $\binom{L_i}{Z_i}$ such sets, together with monotonicity in $S$, gives $V_S=0$ for all
$|S|\ge Z_i$.
\end{itemize}
The excess $(Z_i-D_i)/L_i$ is nondecreasing in $i$ and $1/L_i$ decreases, so $C_n$ attains the
minimum.
\end{proof}

\begin{remark}[Size of the loss]
Writing $Z_i=D_i+e_i$ gives $e_i=2e_{i-1}+g_i$, so $e_i/L_i=e_{i-1}/L_{i-1}+g_i/L_i$. The union bound
needs $g_i\log_2q\gtrsim L_iH_2(Z_i/L_i)+\lambda$. So the relative loss per level is about
$H_2(2^{-k})/\log_2q$, plus $\lambda/(L_i\log_2q)$ on the first few levels. Gates may be sampled in an
extension $\F'\supseteq\F$, at the cost of storing $E_0$ in $\F'$.
\end{remark}

\begin{table}[h]
\centering
\begin{tabular}{lrrrrr}
\toprule
gate field & $n$ & copies & $\Delta$ (Thm.~\ref{thm:distance2}) & queries & $\Delta$ (Thm.~\ref{thm:distance})\\
\midrule
Goldilocks ($2^{64}$) & 20 & 16 & 0.579 & 327 & $<0$\\
Goldilocks ($2^{64}$) & 20 & 32 & 0.729 & 218 & 0.221\\
Goldilocks ($2^{64}$) & 20 & 64 & 0.811 & 183 & 0.605\\
Goldilocks$^3$ ($2^{192}$) & 20 & 4 & 0.675 & 253 & $<0$\\
Goldilocks$^3$ ($2^{192}$) & 20 & 8 & 0.812 & 183 & $<0$\\
M31 ($2^{31}$) & 20 & 64 & 0.417 & 523 & 0.508\\
\bottomrule
\end{tabular}
\caption{Certified distance with setup failure $\le2^{-128}$, and the resulting queries for
$2^{-128}$ soundness (Theorem~\ref{thm:sound}, $\delta=0.95\,\delta_{\max}$). Both theorems are
valid, so the larger value applies.}
\label{tab:distance2}
\end{table}

\begin{theorem}[The generic code is MDS]\label{thm:generic}
Regard the structural points $t^b_i(\omega')$ as independent indeterminates, and work over the
field of rational functions in them. Then every set of at most $D_i$ coordinate functionals of
$C_i$ is linearly independent; equivalently, the coordinate matroid of $C_i$ is uniform of rank
$D_i$, and the generic code $C_n$ is MDS.
\end{theorem}
\begin{proof}
Induction on $i$. For $i=0$ every coordinate functional is $c\mapsto c\neq0$.

\emph{Step.} Identify $\ML_i$ with $\ML_{i-1}^2$ via $F=u+xv$. Position $(\omega',b)$ has functional
$(\ell_{\omega'},\,t^b\ell_{\omega'})$, where $\ell_{\omega'}$ is the coordinate functional of $C_{i-1}$
at $\omega'$. Let $S$ be a set of positions with $|S|\le 2D'$, $D'=D_{i-1}$. Let $A$ be its double
blocks and $B$ its single blocks, with $a=|A|$ and $b=|B|$, so $2a+b\le2D'$.

\emph{Double blocks.} Since $t^1-t^0\neq0$, the two rows of a double block span
$(\ell_{\omega'},0)$ and $(0,\ell_{\omega'})$. By induction the $\ell_{\omega'}$, $\omega'\in A$, are
independent, because $a\le D'$. So the double blocks contribute rank $2a$ and span $U\oplus U$,
with $U=\operatorname{span}\{\ell_{\omega'}:\omega'\in A\}$.

\emph{Quotient.} Pass to $W=\ML_{i-1}^*/U$, of dimension $m=D'-a$, and write
$\bar\ell_{\omega'}$ for the images. By induction, any at most $m$ of the $\bar\ell_{\omega'}$,
$\omega'\in B$, are independent: if $|T|\le m$ then $|A\cup T|\le D'$.

\emph{Single blocks.} Since $b\le2m$, split $B=B_0\sqcup B_1$ with $|B_0|,|B_1|\le m$. Choose
coordinates $J_0$ and $J_1$ of $W$ with $|J_r|=|B_r|$ such that the minors $\det M_{B_0}$ and
$\det M_{B_1}$ of $\{\bar\ell_{\omega'}\}$ are nonzero.

Consider the $b\times b$ minor of the vectors $(\bar\ell_{\omega'},t_{\omega'}\bar\ell_{\omega'})$ on
the columns $J_0$ (first copy of $W$) and $J_1$ (second copy). By the Laplace expansion along the
column split $(J_0,J_1)$, the monomial $\prod_{\omega'\in B_1}t_{\omega'}$ comes only from the term
that puts the rows $B_1$ on $J_1$. Its coefficient is $\pm\det M_{B_0}\det M_{B_1}\neq0$, and it
does not involve the level-$i$ variables. So the minor is nonzero and the $b$ vectors are
independent modulo $U\oplus U$.

The total rank is $2a+b=|S|$.
\end{proof}

\begin{corollary}[Very large fields]\label{cor:generic}
If $q-1\ge2^{\lambda}\,nD\binom{L}{D}$ (with $D=D_n$, $L=L_n$), then with probability at least
$1-2^{-\lambda}$ over uniform gates, $C_n$ is MDS: $\Delta(C_n)=1-(D-1)/L$.
\end{corollary}
\begin{proof}
Parametrise each pair as $(t^0,t^1)=(a,a+s)$ with $(a,s)\in\F\times\F^\times$. This is a
bijection onto distinct pairs, and an invertible linear change of variables.

For $|S|=D$, Theorem~\ref{thm:generic} makes $\det M_S$ a nonzero polynomial. Each entry of $M_S$
is a product of at most $n$ structural points, so $\deg\det M_S\le nD$. Schwartz--Zippel over
$\F\times\F^\times$ gives $\Pr[\det M_S=0]\le nD/(q-1)$. Take a union bound over the
$\binom{L}{D}$ sets.
\end{proof}

\begin{remark}
Corollary~\ref{cor:generic} needs $\log_2q\gtrsim L\,H_2(R)$, so it is structural rather than
practical. It shows that the level count in Theorem~\ref{thm:distance2} is not forced by the
construction. Over Goldilocks with random gates, every $D$-subset was checked to be nonsingular
for $(n,k)\in\{(2,1),(2,2),(2,3),(3,1),(3,2)\}$ (about $1.06\cdot10^{7}$ subsets).

A per-set bound cannot close the gap: for $|S|=D+1$ the measured
$\Pr[\operatorname{rank}M_S<D]$ is about $1.5/q$, not about $q^{-2}$, because a single
coincidence of gate points makes many sets degenerate at once.
\end{remark}

\subsection{Checked setup}

The losses of Theorem~\ref{thm:distance2} at the first levels come from short codes, where a union
bound is expensive. These levels are small enough to be verified once, at setup.

\begin{theorem}[Distance with a checked setup]\label{thm:checked}
Fix $i_0\ge0$. Sample the gates of levels $\le i_0$ (variables $x_n,\dots,x_{n-i_0+1}$) and resample
them until $C_{i_0}$ is MDS. Sample the gates of levels $>i_0$ independently and uniformly on $\Gs(\F)$.
Choose $g_{i_0+1},\dots,g_n\ge0$, put $Z_{i_0}=D_{i_0}$ and $Z_i=2Z_{i-1}+g_i$, and assume $Z_i\le L_i$.
With probability at least
\[
1-\sum_{i=i_0+1}^{n}\binom{L_i}{Z_i}\frac{(1+1/q)^{L_{i-1}}}{q-1}\,q^{-g_i},
\]
every nonzero codeword of $C_n$ has fewer than $Z_n$ zeros, so $\Delta(C_n)\ge1-(Z_n-1)/L_n$.
\end{theorem}
\begin{proof}
$C_{i_0}$ is MDS, so every nonzero codeword of $C_{i_0}$ has at most $D_{i_0}-1$ zeros: $V_S=0$
whenever $|S|\ge Z_{i_0}$. The acceptance test depends only on the gates of levels $\le i_0$, so the
gates of levels $>i_0$ are still independent and uniform after conditioning on acceptance. Run the
induction of Theorem~\ref{thm:distance2} from $i=i_0+1$.
\end{proof}

\begin{remark}[Cost of the check]
$C_1$ is MDS iff the $2^{k+1}$ points of the level-$1$ gates are pairwise distinct. $C_2$ is MDS iff
every $4$-subset of its $2^{k+2}$ positions gives a nonsingular $4\times4$ matrix with rows
$(1,y_1,y_2,y_1y_2)$. By Theorem~\ref{thm:generic} and Schwartz--Zippel, a uniform family is
rejected with probability at most $\binom{L_{i_0}}{D_{i_0}}\,i_0D_{i_0}/(q-1)$, which is negligible for
$i_0\le2$ at $q\approx2^{64}$. The check is done once; the verifier only recomputes the gates from
the public seed and the accepted counter.
\end{remark}

\begin{table}[h]
\centering
\begin{tabular}{lrrrrr}
\toprule
gate field & copies & $i_0=0$ & $i_0=1$ & $i_0=2$ & $i_0=3$\\
\midrule
M31 ($2^{31}$) & 64 & 0.417 & 0.527 & 0.617 & 0.683\\
Goldilocks ($2^{64}$) & 16 & 0.579 & 0.669 & 0.742 & 0.768\\
Goldilocks ($2^{64}$) & 32 & 0.729 & 0.789 & 0.822 & 0.854\\
Goldilocks ($2^{64}$) & 64 & 0.811 & 0.865 & 0.888 & 0.905\\
\bottomrule
\end{tabular}
\caption{Certified $\Delta(C_n)$ from Theorem~\ref{thm:checked}, $n=20$, setup failure $\le2^{-128}$.
The column $i_0=0$ is Theorem~\ref{thm:distance2}.}
\label{tab:checked}
\end{table}

\begin{table}[h]
\centering
\begin{tabular}{rrrrr}
\toprule
$\log_2 q$ (gate field) & copies & $\Delta$ & $1-R$ & queries\\
\midrule
64 & 16 & 0.742 & 0.938 & 210\\
64 & 32 & 0.822 & 0.969 & 180\\
128 & 8 & 0.757 & 0.875 & 201\\
128 & 16 & 0.863 & 0.938 & 169\\
256 & 8 & 0.836 & 0.875 & 176\\
256 & 16 & 0.913 & 0.938 & 157\\
\bottomrule
\end{tabular}
\caption{Gates in a larger field (e.g.\ an extension of the coefficient field), $i_0=2$, $n=20$,
$\lambda=128$. Queries for $2^{-128}$ soundness by Theorem~\ref{thm:sound} with
$\delta=0.95\,\delta_{\max}$. A larger gate field reduces the number of copies; the number of queries
is bounded below by the radius $\delta<\Delta/2$ of Lemma~\ref{lem:agree}.}
\label{tab:extgates}
\end{table}

\subsection{Zero structure}

For $F\in\ML(x_1,\dots,x_n)$ and a prefix $c=(\beta,\theta_n,\dots,\theta_{n-i+1})\in\Omega_i$, let
$F_c\in\ML(x_1,\dots,x_{n-i})$ be $F$ with $x_l$ replaced by $r_l(c)$ for $l>n-i$
(Proposition~\ref{prop:closed}). The prefix $c$ is \emph{killed} if $F_c=0$, and \emph{maximally
killed} if moreover its parent is not killed.

\begin{lemma}[Zero decomposition]\label{lem:zerodec}
Let $F\neq0$ and let $\kappa_i(F)$ be the number of maximally killed prefixes in $\Omega_i$. The
zero set of $\Enc_n(F)$ is the disjoint union of the subtrees below the maximally killed prefixes,
and
\[
\#\{\omega:\Enc_n(F)(\omega)=0\}=\sum_{i=1}^{n}\kappa_i(F)\,2^{n-i}.
\]
Moreover, if $p\in\Omega_{i-1}$ is not killed, at most one child of $p$ is killed. For fixed $F$
and fixed gates of the levels above $p$, a given child of $p$ is killed with probability at most
$1/q$ over the gate at $p$.
\end{lemma}
\begin{proof}
By Proposition~\ref{prop:closed}, $\Enc_n(F)(\omega)=F_\omega$, so the zeros are the killed leaves.
On the path from the root to a zero leaf, $F_{(\beta)}=F\neq0$ and the leaf is killed, so there is a
first killed prefix; it is maximally killed. Every prefix below a killed prefix is killed, since a
restriction of $0$ is $0$. Two maximally killed prefixes are not comparable, so their subtrees are
disjoint, and a prefix in $\Omega_i$ has $2^{n-i}$ leaves.

Write $F_p=A+xB$ with $x=x_{n-i+1}$. The children are $A+t^0(p)B$ and $A+t^1(p)B$. If both vanish,
then $(t^0(p)-t^1(p))B=0$, so $B=0$, $A=0$ and $F_p=0$. A child is killed iff $A=-t\,B$ for its point
$t$: this is impossible if $B=0$ (then $A\neq0$), and otherwise holds for at most one value of $t$.
\end{proof}

\begin{proposition}[Coincidences cost $R\,2^{-j}$]\label{prop:coincidence}
Let $1\le j\le n$ and let $p\neq p'\in\Omega_{j-1}$ lie in different copies ($\beta\neq\beta'$). If
two structural points of their gates coincide, $t^b(p)=t^{b'}(p')=s$, then
\[
\Delta(C_n)\ \le\ 1-R\,(1+2^{-j})+\frac1{L_n}.
\]
\end{proposition}
\begin{proof}
Write $p=(\beta,\theta_n,\dots,\theta_{n-j+2})$. For $1\le i\le j-1$ let $c_i$ be the structural point
of the sibling, off the path of $p$, at level $i$: $c_i=t^{1-\theta_{n-i+1}}_{n-i+1}(\beta,\theta_n,\dots,\theta_{n-i+2})$. Put
\[
F=\prod_{i=1}^{j-1}\bigl(x_{n-i+1}-c_i\bigr)\cdot\bigl(x_{n-j+1}-s\bigr)\cdot K(x_1,\dots,x_{n-j}),
\]
which is multilinear. It vanishes on:
\begin{itemize}[leftmargin=*]
\item the $j-1$ sibling subtrees off the path of $p$, of sizes $2^{n-1},\dots,2^{n-j+1}$, which sum to $2^n-2^{n-j+1}$;
\item the subtrees below $(p,b)$ and $(p',b')$, of size $2^{n-j}$ each. They are disjoint from the
previous ones ($(p,b)$ is on the path of $p$, and $p'$ is in another copy).
\end{itemize}
This gives $2^n$ zeros. The factor $K$ ranges over a space of dimension $2^{n-j}$, so it can be chosen
nonzero and vanishing at $2^{n-j}-1$ further positions. Then $F\neq0$ has at least
$2^n+2^{n-j}-1=D_n-1+2^{n-j}$ zeros, a relative weight of at most $1-R-R\,2^{-j}+1/L_n$.
\end{proof}

\begin{remark}[Consequences for the level count]
Each pair of points of distinct level-$j$ nodes coincides with probability $1/q$. So a theorem with
failure probability $2^{-\lambda}<1/q$ must allow a loss of order $(\lambda/\log_2q)\,R\,2^{-i_0}$ after a
checked setup at level $i_0$. Coincidences are not the only source of excess zeros: for $n=2$, four
positions in different level-$2$ nodes whose points lie on a line $x_2=a+bx_1$ give the nonzero
$F=x_2-a-bx_1$ with $4=D_2$ zeros, although all points are distinct. A proof that removes the
level count must therefore control algebraic relations among many gate points at once.
\end{remark}

\begin{problem}[Removing the level count]\label{prob:levelcount}
Theorems~\ref{thm:distance2} and~\ref{thm:checked} lose about $H_2(2^{-k})/\log_2q$ per level. Prove a bound
\[
\Delta(C_n)\ \ge\ 1-R-c_1\frac{H_2(R)}{\log_2q}-c_2\frac{\lambda}{\log_2q}\,R\,2^{-i_0}-o(1)
\]
with $c_1,c_2$ independent of $n$, or show that a loss growing with $n$ is unavoidable.
\end{problem}

%=====================================================================
\section{Folding}
%=====================================================================

\begin{definition}[Fold]
For $1\le j\le n$ and $z\in\K$, define $\fold_{j,z}:\K^{\Omega_{n-j+1}}\to\K^{\Omega_{n-j}}$ by
\[
(\fold_{j,z}w)(\omega')=w(\omega',0)+\lambda_{j,\omega'}(z)\bigl(w(\omega',1)-w(\omega',0)\bigr),
\qquad
\lambda_{j,\omega'}(z)=\frac{z-t_j^0(\omega')}{t_j^1(\omega')-t_j^0(\omega')}.
\]
By Lemma~\ref{lem:gate-rule}, this equals the gate rule of $A_j(\omega')$ at $x_j=z$:
\[
(\fold_{j,z}w)(\omega')=W_L(z)\,\tfrac{d}{\delta}\,w(\omega',1)-W_R(z)\,\tfrac{c}{\delta}\,w(\omega',0).
\]
\end{definition}

\begin{lemma}[Compatibility]\label{lem:compat}
For $F\in\ML(x_j,\dots,x_n)$,
\[
\fold_{j,z}\bigl(\Enc_{n-j+1}(F)\bigr)=\Enc_{n-j}\bigl(F|_{x_j=z}\bigr).
\]
Consequently, folding $\Enc_n(F)$ with $z_1,\dots,z_n$ gives the constant word $F(z_1,\dots,z_n)$.
\end{lemma}
\begin{proof}
With $F=u+x_jv$, the block at $\omega'$ is $(e_u+t^0e_v,\ e_u+t^1e_v)$, where
$e_u=\Enc(u)(\omega')$ and $e_v=\Enc(v)(\omega')$. Its fold is $e_u+z\,e_v$, and
$\Enc(u)+z\Enc(v)=\Enc(u+zv)$ by linearity.
\end{proof}

\begin{lemma}[Contraction]\label{lem:contract}
$\dpos(\fold_{j,z}w,\fold_{j,z}w')\le\dblk(w,w')$.
\end{lemma}
\begin{proof}
$(\fold w)(\omega')$ depends only on the block of $w$ at $\omega'$.
\end{proof}

\begin{lemma}[Line form]\label{lem:line}
For $w\in\K^{\Omega_{n-j+1}}$ define $A,B\in\K^{\Omega_{n-j}}$ by gate inversion:
\[
B(\omega')=\frac{w(\omega',1)-w(\omega',0)}{t_j^1(\omega')-t_j^0(\omega')},\qquad
A(\omega')=w(\omega',0)-t_j^0(\omega')B(\omega').
\]
Then $\fold_{j,z}w=A+zB$ for all $z$. The map $w\mapsto(A,B)$ is bijective. For $a,b\in C_{n-j}$,
the word $w$ agrees with the codeword of $C_{n-j+1}$ built from $(a,b)$ on exactly the blocks
$\omega'$ where $(A,B)(\omega')=(a,b)(\omega')$.
\end{lemma}
\begin{proof}
Direct from Definition~\ref{def:enc} and $t^0\neq t^1$.
\end{proof}

\begin{lemma}[Line proximity]\label{lem:lineprox}
Let $C\subseteq\K^N$ be linear with $\Delta(C)\ge\Delta$, let $0\le\delta<\Delta/3$, and
$A,B\in\K^N$. Let $\mathcal{Z}=\{z\in\K:\dpos(A+zB,C)\le\delta\}$. If $|\mathcal{Z}|=M\ge2$,
there exist $a,b\in C$ with
\[
\bigl|\{p:(A(p),B(p))\neq(a(p),b(p))\}\bigr|\ \le\ \frac{M}{M-1}\,\delta N.
\]
\end{lemma}
\begin{proof}
Take $z_1\neq z_2$ in $\mathcal{Z}$, with codewords $c_1,c_2$ agreeing with $A+z_iB$ on sets
$S_1,S_2$ of size $\ge(1-\delta)N$. Put $b=(c_1-c_2)/(z_1-z_2)$ and $a=c_1-z_1b$, both in $C$.
On $S_1\cap S_2$ we have $A=a$ and $B=b$.

For any $z_3\in\mathcal{Z}$ with codeword $c_3$ and agreement set $S_3$, the codewords $c_3$ and
$a+z_3b$ agree on $S_1\cap S_2\cap S_3$. This set has size $\ge(1-3\delta)N>(1-\Delta)N$, so
$c_3=a+z_3b$. Hence every $z\in\mathcal{Z}$ has agreement set
$S_z\subseteq\{p:A(p)+zB(p)=a(p)+zb(p)\}$.

Let $E=\{p:(A,B)(p)\neq(a,b)(p)\}$. For $p\in E$ the affine function
$z\mapsto(A-a)(p)+z(B-b)(p)$ is nonzero, so $p\in S_z$ for at most one $z$. Then
$\sum_{z\in\mathcal{Z}}|E\cap S_z|\le|E|$. Combined with $|E\setminus S_z|\le\delta N$, this gives
$M|E|-|E|\le M\delta N$.
\end{proof}

\begin{corollary}[Folding preserves distance]\label{cor:foldprox}
Let $\Delta_{n-j}=\Delta(C_{n-j})$, $0\le\delta<\Delta_{n-j}/3$ and $\eta>0$. If
$\dblk(w,C_{n-j+1}\otimes\K)\ge(1+\eta)\delta$, then
\[
\bigl|\{z\in\K:\ \dpos(\fold_{j,z}w,\ C_{n-j}\otimes\K)\le\delta\}\bigr|\ \le\ 1+\eta^{-1}.
\]
\end{corollary}
\begin{proof}
Write $\fold_{j,z}w=A+zB$ (Lemma~\ref{lem:line}) and let $M$ be the number of such $z$. If $M\ge2$,
Lemma~\ref{lem:lineprox} and Lemma~\ref{lem:line} give $\dblk(w,C_{n-j+1}\otimes\K)\le\frac{M}{M-1}\delta$.
Hence $1+\eta\le M/(M-1)$, i.e.\ $M\le1+\eta^{-1}$. Lemma~\ref{lem:ext} transfers $\Delta$ to $\K$.
\end{proof}

%=====================================================================
\section{The commitment scheme}\label{sec:pcs}
%=====================================================================

\paragraph{Parameters.} $n$ (variables), $k$ (redundancy), $s$ (queries), $\F\subseteq\K$, a hash
function $H$, and a public seed. The gate family $\mathcal{A}$ is derived from the seed with
gates uniform on $\Gs(\F)$. Merkle leaves are blocks: the leaf of a layer at $\omega'$ is the pair
$(w(\omega',0),w(\omega',1))$.

\subsection{Protocol}

\begin{enumerate}[leftmargin=*]
\item \textbf{Commit.} The prover computes $E_0=\Enc_n(F)\in\F^{\Omega_n}$ and publishes its Merkle root $R_0$.
\item \textbf{Fold.} For $j=1,\dots,n$: the verifier sends $z_j\in\K$ uniform (in the non-interactive
version $z_j=H(R_0,\dots,R_{j-1})$). The prover computes $E_j=\fold_{j,z_j}(E_{j-1})\in\K^{\Omega_{n-j}}$.
For $j<n$ it publishes the root $R_j$. For $j=n$ it sends $E_n\in\K^{2^k}$ in full.
\item \textbf{Output.} The verifier checks that $E_n$ is constant and outputs $y=E_n(\beta)$, the
claimed value of $F(z_1,\dots,z_n)$.
\item \textbf{Query.} Repeat $s$ times: sample $\omega=(\beta,\theta_n,\dots,\theta_1)\in\Omega_n$
uniformly. For $j=1,\dots,n$, open the block of $E_{j-1}$ at $\omega^{(j)}$ and the entry
$E_j(\omega^{(j)})$; the latter lies in the block of $E_j$ at $\omega^{(j+1)}$, opened in the next
step, or in $E_n$. Check
\[
E_j(\omega^{(j)})=W_L(z_j)\,\tfrac{d}{\delta}\,E_{j-1}(\omega^{(j)},1)-W_R(z_j)\,\tfrac{c}{\delta}\,E_{j-1}(\omega^{(j)},0)
\]
with $A_j(\omega^{(j)})=\begin{pmatrix}a&b\\c&d\end{pmatrix}$, together with all Merkle paths.
\end{enumerate}

Per query: $n$ block openings, Merkle paths of lengths $n+k-j$, and $O(n)$ field operations.
Prover: $(n+1)2^{n+k}$ operations to encode and $2^{n+k}$ to fold, plus hashing.

%=====================================================================
\section{Security}\label{sec:security}
%=====================================================================

\begin{theorem}[Completeness]
An honest prover is always accepted, and $y=F(z_1,\dots,z_n)$.
\end{theorem}
\begin{proof}
Lemma~\ref{lem:compat} and Lemma~\ref{lem:gate-rule}.
\end{proof}

For words $x,x'$ on the same layer write $\mathrm{Agr}(x,x')=\{\omega:x(\omega)=x'(\omega)\}$, and for
$w,c\in\K^{\Omega_i}$ ($i\ge1$) write $\mathrm{BAgr}(w,c)\subseteq\Omega_{i-1}$ for the set of blocks
on which $w$ and $c$ agree in both entries.

\begin{lemma}[Folding preserves agreement sets]\label{lem:agree}
Let $1\le j\le n$ and $C=C_{n-j}\otimes\K$ with $\Delta(C)\ge\Delta$. Let $\varepsilon>0$ and
\[
0\le\delta<\delta_{\max}(\Delta):=\min\Bigl(1-(1-\Delta)^{1/3},\ \tfrac{\Delta}{2}\Bigr).
\]
Put
\[
\gamma=\frac{(1-\delta)^3-(1-\Delta)}{\Delta},\qquad
M=\Bigl\lfloor\frac{\delta}{\Delta-2\delta}\Bigr\rfloor+2,\qquad
N_1=\Bigl\lceil\max\Bigl(\frac{12}{\Delta\gamma},\ \frac{2M}{\gamma}+2\Bigr)\Bigr\rceil.
\]
For every $w\in\K^{\Omega_{n-j+1}}$ there is a set $B(w)\subseteq\K$ with
$|B(w)|<N_1+1/\varepsilon$ such that, for every $z\notin B(w)$ and every $c'\in C$ with
$\dpos(\fold_{j,z}w,c')\le\delta$, there is $c\in C_{n-j+1}\otimes\K$ with
\[
\fold_{j,z}(c)=c'\qquad\text{and}\qquad
\bigl|\mathrm{Agr}(\fold_{j,z}w,\,c')\setminus\mathrm{BAgr}(w,c)\bigr|\le\varepsilon\,L_{n-j}.
\]
\end{lemma}
\begin{proof}
Write $L=L_{n-j}$ and $\fold_{j,z}w=A+zB$ (Lemma~\ref{lem:line}). Let
$\mathcal{Z}=\{z:\dpos(A+zB,C)\le\delta\}$. Since $2\delta<\Delta$, each $z\in\mathcal{Z}$ has a unique
codeword $c_z$ within distance $\delta$; let $S_z$ be its agreement set, $|S_z|\ge(1-\delta)L$.
Consider the points $P_z=(z,c_z)$ of the affine space $\K\times C$.

\emph{Step 1 (good triples are collinear).} Let $z_1,z_2,z_3\in\mathcal{Z}$ be distinct with
$|S_{z_1}\cap S_{z_2}\cap S_{z_3}|>(1-\Delta)L$. Put $b=(c_{z_1}-c_{z_2})/(z_1-z_2)$ and
$a=c_{z_1}-z_1b$. On $S_{z_1}\cap S_{z_2}$ we have $(A,B)=(a,b)$. On the triple intersection,
$c_{z_3}=A+z_3B=a+z_3b$. Two codewords agreeing on more than $(1-\Delta)L$ positions are equal, so
$c_{z_3}=a+z_3b$: the three points lie on one line.

\emph{Step 2 (many good triples).} Let $N=|\mathcal{Z}|\ge3$ and
$Z_p=\{z\in\mathcal{Z}:p\in S_z\}$. Then $\sum_p|Z_p|\ge N(1-\delta)L$, and the sum of
$|S_{z_1}\cap S_{z_2}\cap S_{z_3}|$ over ordered triples of distinct elements equals
$\sum_p|Z_p|(|Z_p|-1)(|Z_p|-2)$. Since $x(x-1)(x-2)\ge\max(0,x-2)^3$ on integers and the right side
is convex, this sum is at least $L\bigl(N(1-\delta)-2\bigr)^3$. Dividing by $N^3$, the average
triple intersection is at least $L(1-\delta-2/N)^3$. Each triple intersection is at most $L$, and
at most $(1-\Delta)L$ for a bad triple. So the fraction of good triples is at least
\[
\gamma_N=\frac{(1-\delta-2/N)^3-(1-\Delta)}{\Delta}\ \ge\ \gamma-\frac{6}{N\Delta}.
\]

\emph{Step 3 (a heavy line).} Distinct points of $\K\times C$ span a unique line. If line $\ell$
carries $m_\ell$ of the points, ordered collinear triples number
$\sum_\ell m_\ell(m_\ell-1)(m_\ell-2)\le(\max_\ell m_\ell)\,N(N-1)$. Hence some line carries at
least $\gamma_N(N-2)$ points. For $N\ge N_1$ this is at least $M$.

\emph{Step 4 (agreement along a line).} Let a line carry $m\ge2$ points, i.e.\ $c_z=a+zb$ for
$m$ values $z$. Put $E=\{p:(A,B)(p)\neq(a,b)(p)\}$. As in Lemma~\ref{lem:lineprox}, each $p\in E$
lies in at most one $S_z$, while $|E\setminus S_z|\le\delta L$. Hence $|E|\le\frac{m}{m-1}\delta L$.

\emph{Step 5 (absorption).} Let a line $\ell$ with pair $(a,b)$ carry $m\ge M$ points, and let $E$ be as in Step 4, so $|E|\le\frac{M}{M-1}\delta L$. For any $z\in\mathcal{Z}$, the codewords $c_z$ and $a+zb$ both equal $A+zB$ on $S_z\setminus E$, a set of size at least $\bigl(1-\delta-\frac{M}{M-1}\delta\bigr)L$. Since $M-1>\frac{\delta}{\Delta-2\delta}$, we have $\delta\bigl(2+\frac{1}{M-1}\bigr)<\Delta$, so this size exceeds $(1-\Delta)L$ and $c_z=a+zb$. Hence \emph{every} point $P_z$, $z\in\mathcal{Z}$, lies on $\ell$.

\emph{Step 6 (the exceptional set).}
\begin{itemize}[leftmargin=*]
\item If $N<N_1$, put $B(w)=\mathcal{Z}$.
\item Otherwise Step 3 gives a line $\ell$ carrying at least $M$ points, with pair $(a,b)$ and set $E$, and Step 5 puts all of $\mathcal{Z}$ on $\ell$.
\item On $\ell$ we have $S_z=E^{c}\cup(E\cap S_z)$ and $\sum_{z\in\ell}|E\cap S_z|\le|E|\le L$. So
fewer than $1/\varepsilon$ values $z\in\ell$ have $|E\cap S_z|>\varepsilon L$.
\item Put $B(w)=\{z\in\mathcal{Z}:|E\cap S_z|>\varepsilon L\}$.
\end{itemize}
For $z\notin B(w)$ with a codeword $c'$ within $\delta$, we have $z\in\ell$ and $c'=c_z=a+zb$. Let
$c$ be the word of $C_{n-j+1}\otimes\K$ built from $(a,b)$. Then $\fold_{j,z}c=c'$, and
$\mathrm{BAgr}(w,c)=E^{c}$ by Lemma~\ref{lem:line}.
\end{proof}

\begin{remark}
For $\Delta<3-\sqrt5\approx0.76$ the binding term is the triple radius $1-(1-\Delta)^{1/3}$, which
exceeds $\Delta/3$. Example: for $\Delta=0.508$ the radius grows from $0.169$ to $0.21$.
\end{remark}

Fix a family with $\Delta(C_i)\ge\Delta^{*}$ for all $i$, a threshold $0<\delta<\delta_{\max}(\Delta^{*})$ (Lemma~\ref{lem:agree}) and
$\varepsilon>0$. Say the commitment $R_0$ \emph{binds} $F^{*}\in\ML(x_1,\dots,x_n)$ (coefficients in $\K$) if $\dblk(E_0,\Enc_n(F^{*}))<\delta$.
Such an $F^{*}$ is unique, because $2\delta<\Delta^{*}$.

\begin{theorem}[Soundness]\label{thm:sound}
For any prover, the probability that the verifier accepts while either $R_0$ binds no polynomial,
or $R_0$ binds $F^{*}$ and $y\neq F^{*}(z_1,\dots,z_n)$, is at most
\[
\frac{n\,(N_1+1/\varepsilon)}{|\K|}\;+\;\bigl(1-\delta+n\varepsilon\bigr)^{s}
\]
in the IOP model (Merkle binding assumed).
\end{theorem}
\begin{proof}
Write $f_j=\fold_{j,z_j}(E_{j-1})$; the check of step $j$ on a query passes iff
$E_j(\omega^{(j)})=f_j(\omega^{(j)})$.

\emph{Good challenges.} $E_{j-1}$ is fixed before $z_j$ is drawn. Let $G$ be the event
$z_j\notin B(E_{j-1})$ for all $j$ (Lemma~\ref{lem:agree}). Then
$\Pr[\neg G]\le n(N_1+1/\varepsilon)/|\K|$. Assume $G$, and that the verifier checked that $E_n$ is
the constant $y$.

\emph{Descent.} Put $c'_n=E_n\in C_0\otimes\K$. For $j=n,n-1,\dots,1$:
\begin{itemize}[leftmargin=*]
\item if $\dpos(f_j,c'_j)>\delta$, stop and set $J=j$;
\item otherwise, Lemma~\ref{lem:agree} gives $c'_{j-1}$ with $\fold_{j,z_j}c'_{j-1}=c'_j$ and a set
$X_j=\mathrm{Agr}(f_j,c'_j)\setminus\mathrm{BAgr}(E_{j-1},c'_{j-1})$ with
$|X_j|\le\varepsilon L_{n-j}$.
\end{itemize}
If the descent never stops, set $J=0$.
The codewords $c'_j$ and sets $X_j$ depend only on the transcript before the query phase.

\emph{One query.} Let $\omega$ be a query and $p_j=\omega^{(j)}$. Suppose it passes every check
and $p_j\notin X_j$ for all $j>J$. We show by descending induction that $E_j(p_j)=c'_j(p_j)$ for
$j\ge J$.
\begin{itemize}[leftmargin=*]
\item For $j=n$ this is the definition of $c'_n$.
\item Assume it for $j>J$. The check at step $j$ gives $f_j(p_j)=E_j(p_j)=c'_j(p_j)$, so
$p_j\in\mathrm{Agr}(f_j,c'_j)$. Since $p_j\notin X_j$, the block of $E_{j-1}$ at $p_j$ agrees with
$c'_{j-1}$ in both entries. In particular $E_{j-1}(p_{j-1})=c'_{j-1}(p_{j-1})$.
\end{itemize}
Consequently:
\begin{itemize}[leftmargin=*]
\item if $J\ge1$, then $p_J\in\mathrm{Agr}(f_J,c'_J)$, a set of relative size $<1-\delta$;
\item if $J=0$, the block of $E_0$ at $p_1$ lies in $\mathrm{BAgr}(E_0,c'_0)$, a set of relative size
$1-\dblk(E_0,c'_0)$.
\end{itemize}
Since each $p_j$ is uniform on its layer,
\[
\Pr[\text{query passes}]\le n\varepsilon+\max\bigl(1-\delta,\ 1-\dblk(E_0,c'_0)\bigr).
\]

\emph{The bad cases.} Folding $c'_0$ through all rounds gives $c'_n$, the constant $y$. Write
$c'_0=\Enc_n(\widetilde F)$; Lemma~\ref{lem:compat} gives $y=\widetilde F(z)$.
\begin{itemize}[leftmargin=*]
\item If $R_0$ binds no polynomial, then $\dblk(E_0,c'_0)\ge\delta$.
\item If $R_0$ binds $F^{*}$ and $y\neq F^{*}(z)$, then $\widetilde F\neq F^{*}$. Distinct codewords are at
block distance $\ge\Delta^{*}$, so $\dblk(E_0,c'_0)\ge\Delta^{*}-\delta>\delta$.
\end{itemize}
In both cases a query passes with probability $\le1-\delta+n\varepsilon$. The $s$ query positions
are independent of everything fixed above.
\end{proof}

\begin{remark}[Why this bound has no factor $n$]
The argument follows the path of a query from the constant $y$ down to $T_0$. Lemma~\ref{lem:agree}
guarantees that whenever a check passes on a position, the whole pair below it, including the
sibling not on the path, agrees with one fixed valid table. A layer-by-layer distance count
cannot see this, and loses a factor $n$.
\end{remark}

\subsection{Committing only some layers}

Folding is a public deterministic map, so a layer that is not committed can be recomputed by the
verifier from the layer above it. This gives folding by $2^m$ per round with no new soundness loss.

\begin{definition}[Commitment schedule]
A \emph{schedule} is a set $J\subseteq\{1,\dots,n-1\}$ of layers that are \emph{not} committed. The
protocol $\Pi_J$ is the protocol of \S\ref{sec:pcs} with these changes.
\begin{itemize}[leftmargin=*]
\item For $j\in J$ the prover sends no root $R_j$. The challenge $z_{j+1}$ is sampled right after
$z_j$, with no prover message in between (non-interactively: $z_{j+1}=H(\text{transcript},z_j)$).
\item In the query phase, for a committed layer $j_0\notin J$ followed by the run $j_0+1,\dots,j_0+m-1\in J$
and the next committed layer $j_0+m\notin J$ (or $j_0+m=n$), the verifier opens the whole group of
$2^m$ entries of $E_{j_0}$ below the prefix $\omega^{(j_0+m)}$, with one Merkle path. It recomputes the
values of the skipped layers on this group with the gate rule, and checks the resulting value
against the opened entry of $E_{j_0+m}$ (or against $E_n$).
\end{itemize}
Merkle leaves of layer $j_0$ are groups of $2^m$ entries. $J=\emptyset$ is the protocol of
\S\ref{sec:pcs}; $J=\{1,3,5,\dots\}$ is folding by $4$ per round.
\end{definition}

\begin{theorem}[Soundness with any schedule]\label{thm:schedule}
For every $J$, the soundness error of $\Pi_J$ is at most the bound of Theorem~\ref{thm:sound}:
\[
\frac{n\,(N_1+1/\varepsilon)}{|\K|}\;+\;\bigl(1-\delta+n\varepsilon\bigr)^{s}.
\]
\end{theorem}
\begin{proof}
Work in the IOP model of Theorem~\ref{thm:sound}: a root is an oracle to the full table, so the grouping of Merkle leaves plays no role. Let $P^{*}$ be any prover for $\Pi_J$. Build a prover $P'$ for $\Pi_\emptyset$:
\begin{itemize}[leftmargin=*]
\item $P'$ runs $P^{*}$ and forwards its messages and roots.
\item For $j\in J$, after receiving $z_j$, $P'$ defines $E_j:=\fold_{j,z_j}(E_{j-1})$, where $E_{j-1}$ is
the committed table, or the table defined in the same way if $j-1\in J$. It commits to $E_j$ honestly
and sends the root $R_j$.
\item In the query phase $P'$ answers the openings of committed layers with the answers of $P^{*}$,
and the openings of the layers $j\in J$ honestly from $E_j$.
\end{itemize}
\emph{Same challenges.} In both protocols each $z_j$ is uniform on $\K$ and independent of
everything before it. Inserting the root $R_j$ ($j\in J$) does not change this, because $E_j$ and
hence $R_j$ are functions of messages already sent.

\emph{Same acceptance.} Fix a query leaf $\omega$ and a run $j_0,\dots,j_0+m$ as in the definition.
\begin{itemize}[leftmargin=*]
\item The blocks of $E_{j_0}$ that $\Pi_\emptyset$ opens along $\omega$ in this run are contained in the
group that $\Pi_J$ opens, and the $\Pi_J$ verifier sees them under one Merkle path to $R_{j_0}$.
\item For $j\in J$, the $\Pi_\emptyset$ check at layer $j$ compares $E_j(\omega^{(j)})$ with the gate rule
applied to a block of $E_{j-1}$. Since $E_j=\fold_{j,z_j}(E_{j-1})$ entrywise, it passes.
\item The $\Pi_\emptyset$ check at the committed layer $j_0+m$ compares the opened entry of $E_{j_0+m}$
with the gate rule applied to a block of $E_{j_0+m-1}$. That block is the recomputed value of the
group, so this is exactly the $\Pi_J$ check.
\end{itemize}
Hence the $\Pi_\emptyset$ verifier accepts against $P'$ whenever the $\Pi_J$ verifier accepts against
$P^{*}$, on the same randomness, with the same output $y$ and the same root $R_0$. Theorem~\ref{thm:sound}
bounds the probability of the bad event for $P'$.
\end{proof}

\begin{remark}[Costs]
Committing every $m$-th layer gives $\lceil n/m\rceil$ Merkle trees, and per query $\lceil n/m\rceil$
paths, each with a leaf of $2^m$ entries. The verifier work for the gate rule is unchanged, about
$2^m-1$ gate evaluations per committed layer. The same argument covers sending the last $r$ layers
in full: the verifier folds $E_{n-r}$ itself and checks that the result is constant. The query
count $s$ and the challenge field are unchanged.
\end{remark}

\subsection{Batch opening}

To open $t$ polynomials $f_1,\dots,f_t\in\ML(x_1,\dots,x_n)$ at the same random point, the prover
runs one protocol on the $t$-tuple of tables.

\paragraph{Protocol.} Each position carries the tuple $\bigl(E_j^{(1)}(\omega),\dots,E_j^{(t)}(\omega)\bigr)$,
where $E_0^{(i)}=\Enc_n(f_i)$. One Merkle tree per committed layer has these tuples (grouped in blocks) as
leaves. All $t$ tables are folded with the same challenges $z_j$. The last layer is sent for every
$i$; the verifier checks that each component is constant and outputs $y_i=E_n^{(i)}$. A query opens one
path per committed layer and checks the gate rule in every component. The transcript absorbs the whole
last layer before the query positions are derived. No random linear combination of the $f_i$ is used.

For a linear code $C\subseteq\K^L$ let $C^{(t)}=\{(c_1,\dots,c_t):c_i\in C\}$, viewed as a code over the
alphabet $\K^t$: a position is nonzero, or two words agree there, according to the whole $t$-tuple.
Distances, $\mathrm{Agr}$ and $\mathrm{BAgr}$ are taken in this sense.

\begin{proposition}\label{prop:interleave}
$\Delta(C^{(t)})=\Delta(C)$.
\end{proposition}
\begin{proof}
A nonzero $(c_1,\dots,c_t)$ has a nonzero component $c_i$, and its support contains that of $c_i$, so
its weight is at least $\Delta(C)L$. The word $(c^\star,0,\dots,0)$ with $c^\star$ of minimum weight shows equality.
\end{proof}

\begin{lemma}\label{lem:agree-batch}
Lemma~\ref{lem:agree} holds verbatim for $C^{(t)}$ in place of $C$, with the same $\gamma$, $M$, $N_1$ and the
same bound $|B(w)|<N_1+1/\varepsilon$, independent of $t$.
\end{lemma}
\begin{proof}
Folding acts componentwise with the same challenge and gates, so $\fold_{j,z}w=A+zB$ with $A,B$ taking
values in $\K^t$, and the line form (Lemma~\ref{lem:line}) holds componentwise. We check each step of the
proof of Lemma~\ref{lem:agree}.
\begin{itemize}[leftmargin=*]
\item Unique decoding and Steps 1, 5: they use only that $C^{(t)}$ is a $\K$-linear space and that two
codewords agreeing on more than $(1-\Delta)L$ positions are equal, which holds by
Proposition~\ref{prop:interleave}.
\item Steps 2--3: pure counting over positions and over points of the affine space $\K\times C^{(t)}$.
\item Step 4: if $p\in E$ lay in $S_z\cap S_{z'}$, subtracting the two tuple equalities gives
$B(p)=b(p)$ and $A(p)=a(p)$ in $\K^t$, a contradiction; the count is unchanged.
\item Step 6: write $a=(a_i)$, $b=(b_i)$ with $a_i,b_i\in C$ and lift each $a_i+zb_i$ with the line form;
the tuple $c$ of lifts satisfies $\fold_{j,z}c=c_z$ and $\mathrm{BAgr}(w,c)=E^{c}$.\qedhere
\end{itemize}
\end{proof}

\begin{theorem}[Soundness of batch opening]\label{thm:batch}
Say that the commitment binds $(f^\star_1,\dots,f^\star_t)$ if
$\dblk\bigl(E_0,(\Enc_n(f^\star_i))_i\bigr)<\delta$. For any prover, the probability that the verifier
accepts while the commitment binds no tuple, or binds $(f_i^\star)$ but $y_i\neq f_i^\star(z_1,\dots,z_n)$ for
some $i$, is at most
\[
\frac{n\,(N_1+1/\varepsilon)}{|\K|}+\bigl(1-\delta+n\varepsilon\bigr)^{s}.
\]
In particular the number of queries does not depend on $t$. The same holds for every commitment schedule.
\end{theorem}
\begin{proof}
Repeat the proof of Theorem~\ref{thm:sound} with tuple-valued tables, using Lemma~\ref{lem:agree-batch}
once per round (not once per component), so the good-challenge event costs $n(N_1+1/\varepsilon)/|\K|$. A
check passes iff the tuples are equal, so the descent and the one-query induction are unchanged. In the
bad cases, if the descent reaches $J=0$ then $c'_0=(\Enc_n(\tilde f_i))_i$ with $y_i=\tilde f_i(z)$, so
$(\tilde f_i)\neq(f_i^\star)$ and, by uniqueness of the bound tuple (Proposition~\ref{prop:interleave} and
$2\delta<\Delta^{*}$), $\dblk(E_0,c'_0)\ge\delta$. For schedules, the reduction of Theorem~\ref{thm:schedule}
recomputes skipped layers componentwise.
\end{proof}

\begin{remark}[Cost]
The number of queries and Merkle paths is that of a single opening; only the leaves grow, by a factor
$t$ in field elements. The prover encodes and folds $t$ tables.
\end{remark}

\begin{remark}[Non-interactive version]
With $z_j=H(\text{transcript up to }R_{j-1})$ and query positions hashed from the full
transcript, a prover making $Q$ hash queries gains at most a factor $Q$ on the first term, plus the
Merkle collision term.
\end{remark}

\begin{remark}[Concrete parameters, proven]
Take $\F=\mathbb{F}_{2^{31}-1}$, $n=20$, $k=6$ (64 copies, $\Delta^{*}\ge0.508$), $\K$ of degree $5$
over $\F$, $\delta=0.20$ (so $\gamma\approx0.0394$, $M=3$, $N_1=600$) and $\varepsilon=2^{-20}$. Then
$n(N_1+1/\varepsilon)/|\K|\le2^{-130}$, and $s=398$ queries give $(1-\delta+n\varepsilon)^{s}<2^{-128}$.
\end{remark}

\begin{remark}[Any field]
Nothing above uses structure of $\F$ beyond its size: no subgroup and no special element. In particular every result holds in characteristic $2$: the proofs use only $t^0\neq t^1$, division by $t^1-t^0$, and Schwartz--Zippel. In a binary field, additions are XOR and Lemma~\ref{lem:sumrule} reads $\sum_{x\in\{0,1\}}F=c\,F_L+d\,F_R$. If
$\F$ is too small for Theorem~\ref{thm:distance} (roughly $q<2^{30}$ at $n=20$), sample the
gates in an extension $\F'\supseteq\F$. The table $E_0$ then lives in $\F'$, and every
statement holds with $q=|\F'|$. The challenge field $\K$ only needs
$n(N_1+1/\varepsilon)/|\K|$ to be negligible.
\end{remark}

\begin{problem}[Up to $\Delta^{*}/2$]
Remove the triple-radius condition in Lemma~\ref{lem:agree}, so that $\delta<\Delta^{*}/2$ suffices, with a
bad set of size $O(L)+O(1/\varepsilon)$. At $n=20$, four copies, $i_0=3$ and $g=16$ this lowers the query
count from $407$ to $353$. Pairs of close challenges already define lines up to radius $\Delta/2$, and the
absorption step (Step~5 of Lemma~\ref{lem:agree}) holds there; what is missing is the existence of a heavy
line, which for general linear codes in this range is an open question.
\end{problem}

\begin{problem}[Correlated agreement up to the Johnson radius]\label{prob:johnson}
Let $\delta<1-\sqrt{1-\Delta}$ and $\fold_rw=A+rB$. Call a pair $(r,c')$ with $c'$ a codeword at distance
$\le\delta$ from $A+rB$ \emph{good} if some $c$ in the upper code satisfies $\fold_rc=c'$ and
$|\mathrm{Agr}(A+rB,c')\setminus\mathrm{BAgr}(w,c)|\le\varepsilon L$. Show that the challenges having a pair that is
not good number at most $c_1L+c_2/\varepsilon$.

\emph{Evidence.} Exhaustive list decoding of SPARK codes over $\mathrm{GF}(32)$, $\mathrm{GF}(64)$ and
$\mathrm{GF}(256)$ (length $16$, lists of size up to $3$, $1000$ adversarial and planted trials per code and radius)
found at most $2$ such challenges for $\varepsilon L=1$ and at most $1$ for $\varepsilon L=2$, far below the
allowance $1/\varepsilon$. Combined with Theorem~\ref{thm:sound}, a proof would lower the query count at the
parameters above to about $300$.
\end{problem}

\begin{remark}[Structure of the minors of $C_2$]
Write a coordinate row of $C_2$ as $(1,y_1,y_2,y_1y_2)$, with $y_1$ the point of the first gate and $y_2$
that of the child gate. For one copy the single maximal minor is
$\pm(t^1_\varnothing-t^0_\varnothing)^2(t^1_0-t^0_0)(t^1_1-t^0_1)$. For two copies the $70$ maximal minors fall
into three classes by how the four lower blocks are used: $6$ of type $(2,2)$ and $48$ of type $(2,1,1)$
factor into differences of structural points, while the $16$ of type $(1,1,1,1)$, one position from each
lower block, are a single degree-$4$ relation among $8$ points that does not factor into point differences
(checked over $\mathbb Q$, and over $\mathrm{GF}(2^8)$ against all $28$ candidate differences). Hence
pairwise distinctness of gate points certifies every minor containing a double block, but not the
all-single class, whose size grows like $2^{D_i}$. Certifying $C_4$ or $C_5$ cheaply therefore needs a new
idea; it would raise $\Delta^{*}$ to $0.634$ or $0.643$ at the parameters above.
\end{remark}

%=====================================================================
\section{Distance in practice}\label{sec:practice}
%=====================================================================

Theorem~\ref{thm:distance2} loses about $n\,H_2(2^{-k})/\log_2q$ in total. At realistic field
sizes the experiments below found no codeword below the Singleton limit; whether a loss growing
with $n$ is intrinsic is open (Problem~\ref{prob:levelcount}).

\paragraph{Experiment.} Over $\mathbb{F}_{2^{31}-1}$, with per-node random gates:
\begin{itemize}[leftmargin=*]
\item \textbf{Rank tests.} Test sets of $2^n$ positions for linear dependence of their columns,
both uniformly random sets and unions of whole subtrees, where the tree structure would show a
weakness.
\item \textbf{Zero search.} Choose $2^n-1$ positions, take the nonzero polynomial vanishing there,
and count all its zeros. More than $2^n-1$ zeros would beat the Singleton limit.
\end{itemize}

\begin{table}[h]
\centering
\begin{tabular}{rrrrcccc}
\toprule
$n$ & copies & length & $2^n$ & singular (random) & singular (subtrees) & max zeros & $\Delta$ found\\
\midrule
6 & 2 & 128 & 64 & 0/300 & 0/300 & 63 & 0.508\\
6 & 4 & 256 & 64 & 0/300 & 0/300 & 63 & 0.754\\
6 & 8 & 512 & 64 & 0/300 & 0/300 & 63 & 0.877\\
8 & 2 & 512 & 256 & 0/120 & 0/120 & 255 & 0.502\\
8 & 4 & 1024 & 256 & 0/120 & 0/120 & 255 & 0.751\\
8 & 8 & 2048 & 256 & 0/120 & 0/120 & 255 & 0.876\\
\bottomrule
\end{tabular}
\caption{Every test meets the Singleton limit exactly. Control with per-layer gates ($n=6$):
225 to 300 singular sets out of 300, and codewords of relative weight $2^{-6}$, as
Proposition~\ref{prop:perlayer} predicts.}
\label{tab:m31}
\end{table}

\begin{table}[h]
\centering
\begin{tabular}{lrrrrl}
\toprule
setting & copies & $\Delta$ & queries & size of $E_0$ & status\\
\midrule
gates in Goldilocks & 64 & $0.811$ & 183 & $6.7\cdot10^{7}$ & proven (Thm.~\ref{thm:distance2})\\
gates in Goldilocks & 16 & $0.579$ & 327 & $1.7\cdot10^{7}$ & proven (Thm.~\ref{thm:distance2})\\
gates in Goldilocks$^3$ & 8 & $0.812$ & 183 & $8.4\cdot10^{6}$ in $\F_{p^3}$ & proven (Thm.~\ref{thm:distance2})\\
gates in M31 & 64 & $0.508$ & 398 & $6.7\cdot10^{7}$ & proven (Thm.~\ref{thm:distance})\\
any field, $q\ge2^{30}$ & 8 & $\approx0.86$ & $\approx163$ & $8.4\cdot10^{6}$ & if the level count is removed\\
\bottomrule
\end{tabular}
\caption{Parameters for $n=20$ at $2^{-128}$ soundness (Theorem~\ref{thm:sound}).}
\end{table}
%=====================================================================
\section{Implementation and parameters}\label{sec:impl}
%=====================================================================

\paragraph{Instantiation.} Gates in $\F=\mathrm{GF}(2^{128})$ with modulus $x^{128}+x^7+x^2+x+1$ (carry-less
multiplication); challenges and folded layers in $\K=\F[y]/(y^2+y+\alpha)$ with $\mathrm{Tr}(\alpha)=1$, so
$|\K|=2^{256}$; hash BLAKE3; gates derived from a public seed by a keyed hash, one compression per gate, in the
random-oracle model; checked setup at $i_0=3$ (all $\binom{32}{8}$ minors of $C_3$, done once).

\paragraph{Security.} Target $128$ bits against $Q=2^{64}$ hash queries, i.e.\ each term of
Theorem~\ref{thm:sound} at most $2^{-192}$. At $n=20$ and four copies: $\Delta^{*}=0.6157$, $\delta=0.2593$,
$M=4$, $N_1=545$, $\varepsilon=2^{-59.67}$; fold term $192.0$ bits; with $g=16$ bits of grinding, $s=407$ queries
give a query term of $192.3$ bits. These values do not depend on the commitment schedule
(Theorem~\ref{thm:schedule}) nor on the batch size (Theorem~\ref{thm:batch}).

\begin{table}[h]
\centering
\small
\begin{tabular}{rrlrrrrl}
\toprule
$n$ & $s$ & configuration & prover & prover (8 thr.) & verifier & proof & per polynomial\\
\midrule
16 & 376 & single, $m=3$ & 162 ms & 31.7 ms & 6.6 ms & 288 KiB & ---\\
 &  & single, $m=4$ & 124 ms & 24.7 ms & 8.5 ms & 329 KiB & ---\\
 &  & batch $t=16$, $m=3$ & 875 ms & 255 ms & 34.3 ms & 2.98 MiB & 54.7 ms, 2.1 ms, 191 KiB\\
\midrule
18 & 391 & single, $m=3$ & 449 ms & 112 ms & 9.0 ms & 429 KiB & ---\\
 &  & single, $m=4$ & 437 ms & 115 ms & 11.2 ms & 471 KiB & ---\\
 &  & batch $t=16$, $m=3$ & 3.73 s & 1.26 s & 46.0 ms & 3.99 MiB & 233 ms, 2.9 ms, 255 KiB\\
\midrule
20 & 407 & single, $m=3$ & 1.82 s & 0.49 s & 11.4 ms & 580 KiB & ---\\
 &  & single, $m=4$ & 1.72 s & 0.46 s & 14.5 ms & 644 KiB & ---\\
 &  & batch $t=16$, $m=3$ & 17.1 s & 5.87 s & 59.9 ms & 5.26 MiB & 1.07 s, 3.7 ms, 336 KiB\\
\bottomrule
\end{tabular}
\caption{Frozen v0.8 benchmark (tag \texttt{v0.8}, commit \texttt{b8e8081}): four copies, $i_0=3$, $g=16$, Apple M1 (16 GiB), median of 10 runs, every run verified. Prover and verifier single-threaded unless noted; the per-polynomial column gives prover, verifier and proof size divided by $t$. $s$ is derived from Theorem~\ref{thm:sound} at each $n$ for the 192-bit target. Not included: the one-time setup, 0.68 s of gate derivation plus 4.4 s for the $C_3$ check at $n=20$. Peak memory at $n=20$: 0.75 GB for a single opening, 3.2 GB for the batch.}
\label{tab:impl}
\end{table}

%=====================================================================
\section{Further algebraic facts}
%=====================================================================

All statements below were checked numerically over $\mathbb{F}_{2^{31}-1}$ (Rust test suite).

\begin{enumerate}[leftmargin=*]
\item \textbf{Kronecker form.} Let the bottom table $T_0$ come from per-layer gates $A_1,\dots,A_n$,
where layer $i$ splits $x_{n-i+1}$ and children are ordered $(F_L,F_R)$. With the coefficient
vector $\alpha$ indexed so that bit $n-1$ corresponds to $x_n$,
\[
\alpha=(A_1\otimes A_2\otimes\cdots\otimes A_n)\,T_0
\]
($A_1$ acting on the most significant index bit), and
$\det(\bigotimes_iA_i)=\prod_i\det(A_i)^{2^{n-1}}$.
\item \textbf{Sum row.} With $r(A)=[\,2a+c,\ 2b+d\,]$,
$\sum_{x\in\{0,1\}^n}F(x)=\bigl(r(A_1)\otimes\cdots\otimes r(A_n)\bigr)T_0$.
\item \textbf{Representation.} $\rho:\GL_2(\F)^n\to\GL(\F^{2^n})$,
$\rho(A_1,\dots,A_n)=\bigotimes_iA_i$, is a homomorphism with
\[
\ker\rho=\Bigl\{(\lambda_1I,\dots,\lambda_nI):\ \textstyle\prod_i\lambda_i=1\Bigr\}.
\]
The folding map $\alpha\mapsto T_0$ is $\rho(A)^{-1}=\rho(A^{-1})$. The assignment
$A\mapsto\rho(A)^{-1}$ is an anti-homomorphism. The homomorphism to use is
$A\mapsto\rho(A^{-\top})$.
\item \textbf{Identity gate.} For $A_i=I$, $T_0$ is the coefficient vector with index bits
reversed relative to the splitting order.
\end{enumerate}

\end{document}
